% ==================================================================
% §15.5.5 REBUILT — Phase 1: Comparison with other approaches
%
% Part 5 of the §15.5 rebuild. Insert IMMEDIATELY AFTER §15.5.4.
%
% STRICT rules from PROPOSAL v3 Key Change 6:
%   - SHORT categorical comparisons (not a literature survey).
%   - MAXIMUM 1-2 references per category.
%   - Four categories: ΛCDM, EDE, Modified gravity, Phenomenological G_eff.
%
% GPT round 12 corrections (5 edits) applied:
%   (1) Opening: "observational tensions" -> precise listing (Hubble /
%       JWST / late-time consistency channels); retained set isn't
%       all tensions.
%   (2) ΛCDM paragraph: softened; late-time probes noted as part of
%       standard fitting; removed "each observational anomaly addressed
%       through separate" overclaim.
%   (3) EDE: softened "not designed" -> "not primarily constructed /
%       model-dependent".
%   (4) Modified gravity: removed imprecise "typically galactic-scale
%       motivation" (wrong for f(R)); heterogeneous-class framing;
%       added ECT-distinguishing sentence at category level.
%   (5a) Phenomenological G_eff(z): replaced DESI2024 citation with
%       Uzan2011 (proper methodology review).
%   (5b) Key rhetorical point: removed "derived structurally" and
%       "independent observations"; "condensate-closure-motivated
%       drift, treated here at the effective level"; "compatible"
%       instead of "admits".
%
% References used:
%   ΛCDM: Planck2018 (already in bib)
%   EDE: Poulin2019 (added earlier, GPT round 9 / s15.5.5 prep)
%   Modified gravity: Milgrom1983, Sotiriou2010
%   Phenomenological G_eff: Uzan2011 (added 2026-04-18 for GPT r12)
% ==================================================================

\subsubsection*{\normalfont\textbf{15.5.5 \quad Comparison with other approaches}}
\label{subsec:eps_comparison_rebuilt}

The retained five-probe effective band of
\S\ref{subsec:eps_joint_band_rebuilt} can be briefly placed in the
context of alternative proposals addressing the same observational
pattern discussed in
\S\S\ref{subsec:eps_retained_rebuilt}--\ref{subsec:eps_excluded_rebuilt},
namely the Hubble discrepancy, the JWST high-redshift abundance
anomaly, and the associated late-time consistency channels.

\paragraph{$\Lambda$CDM.}
The base $\Lambda$CDM model does not accommodate the Hubble
discrepancy or the JWST early-galaxy excess within its minimal
parameter space. In practice, these tensions are typically discussed
either in terms of residual systematics or by extending the baseline
model in different directions, while the late-time consistency probes
retained in \S\ref{subsec:eps_retained_rebuilt} are already part of
the standard $\Lambda$CDM fitting framework~\cite{Planck2018}.

\paragraph{Early Dark Energy.}
Early-dark-energy proposals primarily target the Hubble discrepancy
through a pre-recombination modification of the expansion history,
typically by introducing one or more additional components or
effective functions~\cite{Poulin2019}. In their basic form they are
not primarily constructed as a simultaneous explanation of the
high-redshift structure-formation anomalies discussed here; whether
they help with the JWST sector is model-dependent and not part of
their minimal motivation.

\paragraph{Modified gravity.}
Modified-gravity approaches form a heterogeneous class. MOND-type
scenarios and TeVeS are historically tied most directly to galactic
phenomenology, while scalar-tensor and $f(R)$-type models are often
motivated by late-time background and growth
modifications~\cite{Milgrom1983,Sotiriou2010}. What distinguishes the
present ECT use of the $\varepsilon$-sector is not simply that gravity
is modified, but that a single effective deformation is confronted
here with both high-redshift and late-time cosmological channels
within one common diagnostic layer.

\paragraph{Phenomenological $G_{\rm eff}(z)$ deformations.}
The $\varepsilon$-deformation used in the present analysis is related
to a broader class of phenomenological varying-coupling or
effective-response parameterizations considered in
cosmology~\cite{Uzan2011}. What distinguishes the ECT use of this
class is that the present effective deformation is motivated by ECT's
three-stage condensate-closure picture and then constrained
empirically here, rather than introduced as an otherwise free
late-time fitting function.

\paragraph{Key rhetorical point.}
The observation that one and the same effective deformation parameter
admits a narrow retained-five-probe effective band across
observationally distinct channels serves as an indirect consistency
check for the framework. One mechanism --- condensate-closure-motivated
drift, treated here at the effective level --- remains compatible
with multiple observational channels within a single narrow band.
This parsimony does not constitute a proof of ECT, but it does
strengthen the plausibility of its cosmological strategy relative to
approaches that require separate parameters or separate mechanisms
for different anomalies.

\paragraph{Honest caveats.}
Three qualifications are stated explicitly. First, the retained
five-probe set is a restricted subset of the broader cosmological
dataset; BAO, $A_{\rm lens}$, and the $S_8$-sector probes contribute
uncertainties and tensions not captured in the joint effective band
(\S\ref{subsec:eps_excluded_rebuilt}). Second, the comparison above is
with schematic representatives of alternative proposals, not with
their full modern pipelines. Third, the ECT treatment of the
$\varepsilon$-sector in \S15.5 is effective, not first-principles;
the closure-derived $\varepsilon(z)$ remains open and is logged as
such in \S15.5.6.

\paragraph{Bridge.}
The implications of the retained five-probe effective band for the
broader ECT programme, and the list of open problems associated with
the $\varepsilon$-sector, are discussed in \S15.5.6.

% ==================================================================
% END §15.5.5 (rebuilt). §15.5.6 will be inserted here in a
% subsequent step. DO NOT REMOVE existing §15.5 below.
% ==================================================================
